% % Neue Befehle
\newcommand{\HRule}[2]{\noindent\rule[#1]{\linewidth}{#2}} % Horiz. Linie
\newcommand{\vlinespace}[1]{\vspace*{#1\baselineskip}} % Abstand
\newcommand{\titleemph}[1]{\textbf{#1}} % Hervorheben

\begin{titlepage}
   \sffamily
   \noindent
   \includegraphics[width=5cm]{fig/aa-titel/HE_Logo_4c}
   \hfill
   \includegraphics[width=5cm]{fig/aa-titel/IPA_Logo.jpg}
      \HRule{13pt}{1pt} 
   \centering
      \Large
      \vlinespace{3}\\
      \workTyp\\
      \huge
      \workTitel\\
%
      \Large
      \vlinespace{2}
      \begin{center}
         in the degree program \\
         {\fontsize{16}{20}\selectfont\workStudiengang} \\
         at the faculty of Mobility and Technology
      \end{center}
%      
      \workSemester\\
%     
      \vlinespace{2}
      \workNameStudent
%
   \vfill
   \raggedright
%   
   \large
   \titleemph{Thesis Period:} \workZeitraum \\ % Nur bei Abschluss-Arbeiten
%   \titleemph{Datum:} \workDatum \\ % Nur bei Studien-Arbeiten
   \titleemph{Supervisor:} Prof. Dr.-Ing. Thao Dang \\
   \titleemph{Second Supervisor:} \workZweitPruefer \\ % Nur bei Abschluss-Arbeiten

 % Folgenden Abschnitt nur bei Industrie-Arbeiten darstellen
   \vlinespace{1}
   \HRule{13pt}{1pt} \\
   \titleemph{Firma:} \workFirma \\
   \titleemph{Betreuer:} \workBetreuer 
%
\end{titlepage}
